% ============================================================
%  يُترجَم بـ XeLaTeX  (في Overleaf: Menu → Compiler → XeLaTeX)
%  الخط: Amiri (متوفر في Overleaf، أو حمّله من github.com/aliftype/amiri)
% ============================================================
\documentclass[12pt,a4paper]{article}
\usepackage[a4paper,margin=1.7cm,top=2.3cm,bottom=2.9cm,headheight=15pt,footskip=1.5cm]{geometry}
\usepackage{amsmath,amssymb}
\usepackage{icomma}
\usepackage[english,bidi=default,layout=lists]{babel}
\babelprovide[import,main]{arabic}
\babelfont{rm}[Scale=1.08]{Amiri}
\babelfont{sf}[Scale=1.08]{Amiri}
\usepackage[most]{tcolorbox}
\usepackage{tikz}
\usetikzlibrary{arrows.meta,calc,shadings}
\usepackage{array}
\usepackage{colortbl}
\usepackage{fancyhdr}
\usepackage{titlesec}
\usepackage{enumitem}
\usepackage{collcell}
\usepackage{needspace}
\usepackage{etoolbox}
\newcolumntype{C}{>{\collectcell\ar}c<{\endcollectcell}}
\pretocmd{\section}{\Needspace{15\baselineskip}}{}{}
\pretocmd{\subsection}{\Needspace{11\baselineskip}}{}{}

\definecolor{cDef}{HTML}{1F4E9C}
\definecolor{cThm}{HTML}{B3261E}
\definecolor{cRem}{HTML}{2E7D32}
\definecolor{cEx}{HTML}{C76B00}
\definecolor{cMeth}{HTML}{6A1B9A}
\definecolor{cSol}{HTML}{00695C}
\definecolor{cSec}{HTML}{1F3A68}

\titleformat{\section}{\Large\bfseries\color{cSec}}{\thesection.}{0.6em}{}[{\color{cSec!60}\titlerule[1pt]}]
\titleformat{\subsection}{\large\bfseries\color{cThm!85!black}}{\thesubsection.}{0.5em}{}
\titlespacing*{\section}{0pt}{14pt}{8pt}
\titlespacing*{\subsection}{0pt}{10pt}{5pt}
\setlength{\parindent}{0pt}
\setlength{\parskip}{4pt}
\setlist{itemsep=2pt,topsep=3pt}

\newtcolorbox{arbox}[3][]{enhanced,
  colback=#2!5,colframe=#2!80!black,arc=2mm,boxrule=0.8pt,
  left=6pt,right=6pt,top=8pt,bottom=5pt,
  fonttitle=\bfseries,coltitle=white,
  title={\foreignlanguage{arabic}{#3}},
  attach boxed title to top right={xshift=-5mm,yshift=-3mm},
  boxed title style={colback=#2!80!black,arc=1.5mm,boxrule=0pt},
  before upper={\selectlanguage{arabic}},
  before skip=8pt,after skip=8pt,breakable,#1}
\newenvironment{BOX}[2]{\begin{otherlanguage}{english}\begin{arbox}{#1}{#2}}{\end{arbox}\end{otherlanguage}}
\newenvironment{tarif}[1][تعريف]{\begin{BOX}{cDef}{#1}}{\end{BOX}}
\newenvironment{khassia}[1][خاصية]{\begin{BOX}{cThm}{#1}}{\end{BOX}}
\newenvironment{molahada}[1][ملاحظة]{\begin{BOX}{cRem}{#1}}{\end{BOX}}
\newenvironment{mithal}[1][مثال]{\begin{BOX}{cEx}{#1}}{\end{BOX}}
\newenvironment{tariqa}[1][طريقة]{\begin{BOX}{cMeth}{#1}}{\end{BOX}}

\newcounter{exo}
\newenvironment{tamrin}[1]{\stepcounter{exo}\begin{BOX}{cSec}{التمرين \theexo: #1}}{\end{BOX}}
\newenvironment{hal}{\Needspace{12\baselineskip}\begin{otherlanguage}{english}\begin{arbox}[breakable]{cSol}{حل التمرين \theexo}}{\end{arbox}\end{otherlanguage}}

\newenvironment{LTR}{\begin{otherlanguage}{english}}{\end{otherlanguage}}
\newcommand{\ar}[1]{\foreignlanguage{arabic}{#1}}

\renewcommand{\dots}{\mathinner{\ldotp\ldotp\ldotp}}
\renewcommand{\ldots}{\mathinner{\ldotp\ldotp\ldotp}}
\renewcommand{\cdots}{\mathinner{\cdotp\cdotp\cdotp}}
\newcommand{\N}{\mathbb{N}}
\newcommand{\Z}{\mathbb{Z}}
\newcommand{\D}{\mathbb{D}}
\newcommand{\Q}{\mathbb{Q}}
\newcommand{\R}{\mathbb{R}}
\newcommand{\pgcd}{\operatorname{PGCD}}
\newcommand{\ppcm}{\operatorname{PPCM}}

% \NLseg{a}{b}{[}{]} ; \NLright{a}{[} ; \NLleft{b}{]}  (interval on a number line)
\newcommand{\NLseg}[4]{\begin{LTR}\begin{tikzpicture}[baseline=-0.6ex]
 \draw[-{Stealth[length=2mm]},gray!70](0,0)--(3.4,0);
 \draw[line width=2pt,cThm!80!black](0.7,0)--(2.7,0);
 \node[cThm!80!black,font=\LARGE] at (0.7,0){$#3$};
 \node[cThm!80!black,font=\LARGE] at (2.7,0){$#4$};
 \node[below=3pt,font=\footnotesize] at (0.7,0){$#1$};
 \node[below=3pt,font=\footnotesize] at (2.7,0){$#2$};
\end{tikzpicture}\end{LTR}}
\newcommand{\NLright}[2]{\begin{LTR}\begin{tikzpicture}[baseline=-0.6ex]
 \draw[-{Stealth[length=2mm]},gray!70](0,0)--(3.4,0);
 \draw[line width=2pt,cThm!80!black](1.3,0)--(3.3,0);
 \node[cThm!80!black,font=\LARGE] at (1.3,0){$#2$};
 \node[below=3pt,font=\footnotesize] at (1.3,0){$#1$};
\end{tikzpicture}\end{LTR}}
\newcommand{\NLleft}[2]{\begin{LTR}\begin{tikzpicture}[baseline=-0.6ex]
 \draw[-{Stealth[length=2mm]},gray!70](0,0)--(3.4,0);
 \draw[line width=2pt,cThm!80!black](0.05,0)--(2.1,0);
 \node[cThm!80!black,font=\LARGE] at (2.1,0){$#2$};
 \node[below=3pt,font=\footnotesize] at (2.1,0){$#1$};
\end{tikzpicture}\end{LTR}}

% ---------- social-media icons (pure TikZ) ----------
\definecolor{cFB}{HTML}{1877F2}
\definecolor{cTG}{HTML}{229ED9}
\definecolor{cIG1}{HTML}{FEDA75}
\definecolor{cIG2}{HTML}{D62976}
\definecolor{cIG3}{HTML}{4F5BD5}
\newcommand{\fbicon}{\begin{LTR}\begin{tikzpicture}[baseline=-0.75ex,x=1em,y=1em]
 \fill[cFB] (0,0) circle (0.55);
 \node[white,inner sep=0,font=\usefont{TU}{lmss}{bx}{n}\small] at (0.04,-0.02) {f};
\end{tikzpicture}\end{LTR}}
\newcommand{\igicon}{\begin{LTR}\begin{tikzpicture}[baseline=-0.75ex,x=1em,y=1em]
 \shade[lower left=cIG1,lower right=cIG2,upper left=cIG2,upper right=cIG3,rounded corners=0.28em] (-0.55,-0.55) rectangle (0.55,0.55);
 \draw[white,line width=0.09em,rounded corners=0.18em] (-0.36,-0.36) rectangle (0.36,0.36);
 \draw[white,line width=0.09em] (0,0) circle (0.17);
 \fill[white] (0.22,0.22) circle (0.05);
\end{tikzpicture}\end{LTR}}
\newcommand{\tgicon}{\begin{LTR}\begin{tikzpicture}[baseline=-0.75ex,x=1em,y=1em]
 \fill[cTG] (0,0) circle (0.55);
 \fill[white] (-0.33,0.0) -- (0.32,0.27) -- (0.2,-0.3) -- (0.0,-0.15) -- (-0.08,-0.28) -- (-0.08,-0.08) -- cycle;
 \fill[cTG!70!black] (-0.08,-0.08) -- (0.25,0.2) -- (0.0,-0.15) -- (-0.08,-0.28) -- cycle;
\end{tikzpicture}\end{LTR}}
% page number: black box, white two-digit number, at the right of the social row
\newcommand{\PageBox}{\begin{LTR}\tikz[baseline=(p.base)]\node[fill=black,text=white,inner sep=2.5pt,minimum width=1.6em,minimum height=1.3em,font=\usefont{TU}{lmss}{bx}{n}\small](p){\ifnum\value{page}<10 0\fi\arabic{page}};\end{LTR}}
% visual left-to-right order: Telegram | Instagram | Facebook | page box
\newcommand{\SocialRow}{\begin{otherlanguage}{english}%
 \ar{\textbf{الأستاذ سوايسية هشام للرياضيات}}\ \tgicon\hfill
 \textbf{prof\_hicham\_math}\ \igicon\hfill
 \ar{\textbf{الأستاذ سوايسية هشام للرياضيات}}\ \fbicon\hspace{2.2em}\PageBox\end{otherlanguage}}

\pagestyle{fancy}
\fancyhf{}
\renewcommand{\headrulewidth}{0.6pt}
\renewcommand{\footrulewidth}{0.4pt}
% \SetHeader{right head: document title}{left head: level -- stream}
\newcommand{\SetHeader}[2]{%
 \fancyhead[C]{\parbox{\textwidth}{\ar{#1\hfill #2}}}%
 \fancyfoot[C]{\parbox{\textwidth}{\small\SocialRow}}}
\usepackage{tkz-tab}
\usepackage{multicol}
% variation / sign table (LTR)
\newenvironment{TAB}{\par\begin{center}\begin{LTR}\begin{tikzpicture}[scale=0.95,every node/.style={font=\small}]}{\end{tikzpicture}\end{LTR}\end{center}}
% grid + axes, clipped: \begin{GR}{xmin}{xmax}{ymin}{ymax}{unit cm}
\newenvironment{GR}[5]{\par\begin{center}\begin{LTR}\begin{tikzpicture}[x=#5cm,y=#5cm,font=\footnotesize]
 \draw[gray!25,very thin] (#1,#3) grid (#2,#4);
 \draw[-{Stealth[length=2mm]}] (#1,0)--(#2+0.3,0) node[below left]{$x$};
 \draw[-{Stealth[length=2mm]}] (0,#3)--(0,#4+0.3) node[below left]{$y$};
 \node[below left] at (0,0){$O$};
 \begin{scope}\clip (#1,#3) rectangle (#2,#4);}
{\end{scope}\end{tikzpicture}\end{LTR}\end{center}}
\newcommand{\Ticks}[2]{\foreach \t in {#1} \node[below,font=\scriptsize] at (\t,0){$\t$}; \foreach \t in {#2} \node[left,font=\scriptsize] at (0,\t){$\t$};}
\newcommand{\bx}[1]{\ensuremath{\colorbox{cDef!12}{$\displaystyle #1$}}}
\newcommand{\TwoCol}[4]{\par\begin{LTR}\noindent\begin{minipage}[t]{#1\textwidth}\selectlanguage{arabic}#2\end{minipage}\hfill\begin{minipage}[t]{#3\textwidth}\selectlanguage{arabic}#4\end{minipage}\end{LTR}\par}
\newcommand{\RL}[4]{\TwoCol{#3}{#4}{#1}{#2}}  % \RL{w1}{right content}{w2}{left content}
\newcommand{\ok}{$\textcolor{cRem}{\checkmark}$}
\newcommand{\limn}{\displaystyle\lim_{n\to+\infty}}
% sub-questions, each on its own line
\newenvironment{sq}{\begin{itemize}[leftmargin=2.4em,itemsep=3pt]}{\end{itemize}}
\newcommand{\qa}{\item[(أ)]}\newcommand{\qb}{\item[(ب)]}\newcommand{\qc}{\item[(جـ)]}\newcommand{\qd}{\item[(د)]}\newcommand{\qe}{\item[(هـ)]}
\newcommand{\pa}{\par\smallskip\textbf{(أ)} }\newcommand{\pb}{\par\smallskip\textbf{(ب)} }\newcommand{\pc}{\par\smallskip\textbf{(جـ)} }\newcommand{\pd}{\par\smallskip\textbf{(د)} }\newcommand{\pe}{\par\smallskip\textbf{(هـ)} }
\newcommand{\QQ}[1]{\par\smallskip\textbf{#1.} }
\newenvironment{tmr}{\stepcounter{exo}\begin{BOX}{cSec}{تمرين محلول \theexo}}{\end{BOX}}
\newcommand{\ABC}[3]{\\[2pt]\begin{tabular}{p{0.28\textwidth}p{0.28\textwidth}p{0.28\textwidth}}(جـ) \ #3 & (ب) \ #2 & (أ) \ #1\end{tabular}}
\pretocmd{\section}{\Needspace{24\baselineskip}}{}{}
\newtcbox{\bxtbox}{on line,colback=cDef!12,colframe=cDef!12,boxrule=0pt,arc=1pt,boxsep=0pt,left=2pt,right=2pt,top=1pt,bottom=1pt}
\newcommand{\bxt}[1]{\foreignlanguage{english}{\bxtbox{\ar{#1}}}}
\usepackage{longtable}
\setlength{\LTpre}{4pt}\setlength{\LTpost}{0pt}
\newcommand{\Phase}[1]{\textbf{#1}}
\newcommand{\Fld}[2]{\textbf{#1:}\ #2}
% رأس المذكرة على نمط النموذج: إطار بعنوان «مذكرة رقم»، معلومات الحصة، ثم إطار الكفاءات
\newcommand{\ModhHead}[9]{%
\begin{LTR}
\begin{tcolorbox}[enhanced,colback=white,colframe=cThm!85!black,arc=4mm,boxrule=1pt,top=14pt,bottom=6pt,left=8pt,right=8pt,
 overlay={\node[fill=cThm!85!black,text=white,rounded corners=3mm,inner xsep=14pt,inner ysep=4pt,font=\large\bfseries] at ([yshift=0pt]frame.north){\ar{مذكرة رقم #1}};},
 before upper={\selectlanguage{arabic}}]
\begin{tabular}{@{}p{0.4\linewidth}p{0.55\linewidth}@{}}
\Fld{المستوى}{#2} & \Fld{الثانوية}{\dotfill}\\[3pt]
\Fld{المادة}{رياضيات} & \Fld{ميدان التعلم}{#3}\\[3pt]
\Fld{المدة}{#5} & \Fld{الوحدة التعلمية}{#4}\\[3pt]
\Fld{الأستاذ}{\dotfill} & \Fld{المحتوى المعرفي}{#6}\\
\end{tabular}
\end{tcolorbox}
\begin{tcolorbox}[enhanced,colback=white,colframe=cThm!85!black,arc=3mm,boxrule=0.8pt,sidebyside,sidebyside align=top,lefthand ratio=0.5,
 segmentation style={dashed,cThm!60},before upper={\selectlanguage{arabic}},before lower={\selectlanguage{arabic}}]
\textbf{الكفاءات المستهدفة:}\par #8
\tcblower
\textbf{الكفاءات القبلية:}\par #7
\end{tcolorbox}
\end{LTR}
\par\smallskip\textbf{الوسائل:} #9\par\smallskip}
\newenvironment{Sayr}{%
\begin{LTR}\renewcommand{\arraystretch}{1.25}
\begin{longtable}{|p{0.17\textwidth}|>{\centering\arraybackslash}p{0.06\textwidth}|p{0.53\textwidth}|>{\centering\arraybackslash}p{0.12\textwidth}|}
\hline\rowcolor{cThm!12}\ar{\textbf{توجيهات وتعاليق}} & \ar{\textbf{المدة}} & \ar{\textbf{سير الحصة}} & \ar{\textbf{مراحل الحصة}}\\\hline\endhead
\hline\endfoot}{\end{longtable}\end{LTR}}
% سطر في الجدول: \Row{مرحلة}{سير}{مدة}{توجيهات}
\newcommand{\Row}[4]{\ar{\small #4} & \ar{#3} & \ar{#2} & \cellcolor{cThm!8}\ar{\Phase{#1}}\\\hline}

\newcommand{\Lab}[2]{\foreignlanguage{english}{\tcbox[on line,colback=#1,colframe=#1,boxrule=0pt,arc=2pt,boxsep=1pt,left=5pt,right=5pt,top=1pt,bottom=1pt]{\ar{\textcolor{white}{\textbf{#2}}}}}}
\newcommand{\Num}[1]{\foreignlanguage{english}{\tikz[baseline=(c.base)]\node[fill=cDef,text=white,inner sep=2pt,font=\footnotesize\bfseries](c){#1};}}
\newcommand{\ck}{\ensuremath{{\color{cRem}\checkmark}}}
\newenvironment{Kad}[1]{\par\smallskip\Lab{#1}}{\par}
\SetHeader{\small\bfseries\color{cSec}مذكرات محور: المتتاليات العددية}{\small\color{cSec}السنة الثانية ثانوي -- علوم تجريبية}
\begin{document}
\begin{LTR}
\begin{tcolorbox}[enhanced,colback=cThm!85!black,colframe=cThm!85!black,arc=3mm,halign=center,boxrule=0pt,top=12pt,bottom=12pt,before upper={\selectlanguage{arabic}}]
{\color{white}\fontsize{26}{32}\selectfont\bfseries مذكرات محور: المتتاليات العددية}\\[8pt]
{\color{white}\large السنة الثانية ثانوي -- علوم تجريبية (صالحة لشعبتي رياضيات وتقني رياضي)}\\[4pt]
{\color{white}\large 8 مذكرات (16 ساعة) -- السنة الدراسية 2026-2027}
\end{tcolorbox}
\end{LTR}
\begin{BOX}{cThm}{يتضمن هذا العمل}
\begin{enumerate}[label=\textbf{مذكرة \arabic{enumi}:},leftmargin=6em,start=56]
\item عموميات: توليد متتالية عددية.
\item التمثيل البياني لمتتالية.
\item اتجاه تغير متتالية.
\item المتتاليات الحسابية.
\item مجموع حدود متتابعة لمتتالية حسابية ، الوسط الحسابي.
\item المتتاليات الهندسية.
\item مجموع حدود متتابعة لمتتالية هندسية ، الوسط الهندسي.
\item نهاية متتالية ، المتتاليات المتقاربة.
\end{enumerate}
مراجع النشاطات والتمارين: الكتاب المدرسي للسنة الثانية ثانوي (باب المتتاليات ص 178--179). حسب المنهاج: لا برهان بالتراجع ولا متتاليات محدودة في هذا المستوى.
\end{BOX}
\newpage
\ModhHead{56}{2 ع ت}{تحليل}{المتتاليات العددية}{02 سا}{عموميات: توليد متتالية عددية}%
{\begin{itemize}\item الدوال ، العبارات الجبرية.\end{itemize}}%
{\begin{itemize}\item توليد متتالية بدستور صريح أو بعلاقة تراجعية.\item وصف ظاهرة بواسطة متتالية.\end{itemize}}%
{المنهاج ، الكتاب المدرسي ، مراجع ، الحاسبة.}
\begin{Sayr}
\Row{مرحلة الانطلاق}{%
\textbf{التهيئة النفسية:} $0$ ، $5$ ، $10$ ، $15$ ، $\ldots$: ما العدد الموالي؟ وما العدد الذي في المرتبة $100$ ؟}{10 د}{لغز عددي.}
\Row{التشخيص والاكتشاف}{%
\Lab{cEx}{مناقشة النشاطين 01 و 02 صفحة 178}\par\ck\ مضاعفات $5$: $u_n=5n$ (دستور صريح بدلالة $n$).\par
\ck\ عدد المثلثات: للحصول على $u_{n+1}$ نضيف إلى $u_n$ عدداً: علاقة بين حدين متتابعين (علاقة تراجعية).\par
}{35 د}{نشاطا الكتاب.}
\Row{بناء المعارف}{%
\Lab{cDef}{تعريف}\ \textbf{المتتالية العددية} دالة من $\mathbb N$ (أو جزء منه) نحو $\mathbb R$ ؛ نرمز $u_n$ للحد ذي الدليل $n$ ، و $(u_n)$ للمتتالية.\par\Lab{cMeth}{طرق التوليد}\par\ck\ \textbf{صريحة}: $u_n=f(n)$ ، مثلاً $u_n=n^2-1$: $u_0=-1$ ، $u_1=0$ ، $u_2=3$.\par
\ck\ \textbf{تراجعية}: $u_0=2$ و $u_{n+1}=3u_n-1$: $u_1=5$ ، $u_2=14$ ، $u_3=41$.\par
}{50 د}{نميز الطريقتين في جدول.}
\Row{}{%
\Lab{cEx}{وصف ظاهرة}\ مدينة سكانها $P_0=1000$ يزدادون بـ $5\,\%$ كل سنة: $P_{n+1}=1{,}05\,P_n$ ؛ $P_1=1050$ ، $P_2=1102{,}5$.\par\Lab{cThm!80!black}{حذار}\ في العلاقة التراجعية لا نستطيع حساب $u_{10}$ دون الحدود السابقة.}{}{}
\Row{التقويم}{%
\Lab{cEx}{تطبيق}\par احسب الحدود الثلاثة الأولى: \Num{1} $u_n=\dfrac{2n+1}{n+1}$ ؛ \Num{2} $v_0=1$ ، $v_{n+1}=v_n+2n$.\par\textbf{الحل:} \Num{1} $1$ ، $\frac32$ ، $\frac53$ ؛ \Num{2} $v_1=1$ ، $v_2=3$ ، $v_3=7$.\par\textbf{واجب منزلي:} تمارين 20 ، 21 ، 24 ، 27 ، 30 من باب المتتاليات.\par\smallskip\textbf{ملاحظات حول سير الحصة:} \dotfill}{25 د}{تقويم فردي.}
\end{Sayr}
\newpage
\ModhHead{57}{2 ع ت}{تحليل}{المتتاليات العددية}{02 سا}{التمثيل البياني لمتتالية}%
{\begin{itemize}\item توليد متتالية ، التمثيل البياني لدالة.\end{itemize}}%
{\begin{itemize}\item تمثيل حدود متتالية بيانياً ، وتخمين سلوكها.\end{itemize}}%
{المنهاج ، الكتاب المدرسي ، مراجع ، الحاسبة. ورق ميليمتري.}
\begin{Sayr}
\Row{مرحلة الانطلاق}{%
\textbf{التهيئة النفسية:} كيف نمثل المتتالية $u_n=n^2$ في معلم؟ هل نصل النقط؟}{10 د}{سؤال.}
\Row{التشخيص والاكتشاف}{%
\Lab{cEx}{مناقشة نشاط 04 صفحة 179}\par\ck\ دالة $f$ ومستقيم $y=x$: نقرأ $u_1=f(u_0)$ على المنحني ، ثم ننقله إلى محور الفواصل بالمستقيم $y=x$ ، وهكذا.\par
}{35 د}{نشاط الكتاب.}
\Row{بناء المعارف}{%
\Lab{cMeth}{طريقة}\par\ck\ $u_n=f(n)$: نمثل النقط المنعزلة $(n;u_n)$ ولا نصلها.\par
\ck\ $u_{n+1}=f(u_n)$: نرسم $C_f$ والمستقيم $y=x$ ، ونمثل الحدود على محور الفواصل بطريقة «السلم».\par
\Lab{cEx}{مثال}\ $f(x)=\frac12x+1$ ، $u_0=0$: $u_1=1$ ، $u_2=1{,}5$ ، $u_3=1{,}75$.\par\smallskip\begin{LTR}\centering \begin{tikzpicture}[scale=2.4,font=\scriptsize]
\draw[->] (-0.2,0)--(2.6,0) node[right]{$x$}; \draw[->] (0,-0.2)--(0,2.4) node[above]{$y$};
\draw[cRem,dashed] (0,0)--(2.3,2.3) node[above]{$y=x$};
\draw[cDef,thick] (-0.1,0.95)--(2.5,2.25) node[right]{$y=f(x)$};
\draw[cThm,thick,->] (0,0)--(0,1)--(1,1)--(1,1.5)--(1.5,1.5)--(1.5,1.75)--(1.75,1.75);
\foreach \x/\n in {0/0,1/1,1.5/2,1.75/3} {\draw[dotted] (\x,0)--(\x,\x); \node[below] at (\x,0){$u_{\n}$};}
\end{tikzpicture}\end{LTR}\par }{50 د}{رسم السلم بلونين.}
\Row{}{%
\Lab{cRem}{ملاحظات}\ نخمن أن الحدود تقترب من $2$ (فاصلة نقطة تقاطع $C_f$ مع $y=x$).\par\Lab{cThm!80!black}{حذار}\ التمثيل يسمح بالتخمين فقط ؛ البرهان جبري.}{}{}
\Row{التقويم}{%
\Lab{cEx}{تطبيق}\par $u_0=4$ و $u_{n+1}=\frac12u_n+1$. مثل الحدود $u_0$ إلى $u_3$ بيانياً ، وخمن سلوك المتتالية.\par\textbf{الحل:} $u_1=3$ ، $u_2=2{,}5$ ، $u_3=2{,}25$: متناقصة وتقترب من $2$.\par\textbf{واجب منزلي:} تمارين من باب المتتاليات.\par\smallskip\textbf{ملاحظات حول سير الحصة:} \dotfill}{25 د}{تقويم فردي.}
\end{Sayr}
\newpage
\ModhHead{58}{2 ع ت}{تحليل}{المتتاليات العددية}{02 سا}{اتجاه تغير متتالية}%
{\begin{itemize}\item توليد متتالية ، اتجاه تغير دالة.\end{itemize}}%
{\begin{itemize}\item دراسة اتجاه تغير متتالية ، أو ابتداءً من رتبة معينة.\end{itemize}}%
{المنهاج ، الكتاب المدرسي ، مراجع ، الحاسبة.}
\begin{Sayr}
\Row{مرحلة الانطلاق}{%
\textbf{التهيئة النفسية:} $u_n=n^2-5n$: احسب $u_0$ إلى $u_5$. هل الحدود تتزايد دائماً؟}{10 د}{حساب حدود.}
\Row{التشخيص والاكتشاف}{%
\Lab{cEx}{وضعية مقترحة}\par\ck\ $0$ ، $-4$ ، $-6$ ، $-6$ ، $-4$ ، $0$: تتناقص ثم تتزايد.\par
\ck\ $u_{n+1}-u_n=2n-4\geqslant0$ لما $n\geqslant2$: متزايدة \textbf{ابتداءً من الرتبة} $2$.\par
}{35 د}{جدول قيم.}
\Row{بناء المعارف}{%
\Lab{cDef}{تعريف}\ $(u_n)$ متزايدة (ابتداءً من رتبة $n_0$) إذا كان $u_{n+1}\geqslant u_n$ من أجل كل $n\geqslant n_0$ ؛ ومتناقصة إذا كان $u_{n+1}\leqslant u_n$.\par\Lab{cMeth}{طريقة}\par\ck\ ندرس إشارة الفرق $u_{n+1}-u_n$.\par
\ck\ إذا كان $u_n=f(n)$ و $f$ رتيبة على $[0;+\infty[$ ، فلـ $(u_n)$ نفس اتجاه تغير $f$.\par
\ck\ إذا كانت الحدود موجبة تماماً: نقارن $\dfrac{u_{n+1}}{u_n}$ بـ $1$.\par
}{50 د}{الطرق الثلاث في جدول.}
\Row{}{%
\Lab{cEx}{مثال}\ $u_n=\dfrac1{n+1}$: $f(x)=\dfrac1{x+1}$ متناقصة على $[0;+\infty[$ ، إذن $(u_n)$ متناقصة.\par\Lab{cThm!80!black}{حذار}\ عكس الخاصية الثانية غير صحيح عموماً.}{}{}
\Row{التقويم}{%
\Lab{cEx}{تطبيق}\par ادرس اتجاه تغير $v_n=n^2-6n+1$.\par\textbf{الحل:} $v_{n+1}-v_n=2n-5>0$ لما $n\geqslant3$: $(v_n)$ متزايدة ابتداءً من الرتبة $3$.\par\textbf{واجب منزلي:} تمارين 34 ، 35 ، 36 ، 45 ، 46 ، 47 من باب المتتاليات.\par\smallskip\textbf{ملاحظات حول سير الحصة:} \dotfill}{25 د}{تقويم فردي.}
\end{Sayr}
\newpage
\ModhHead{59}{2 ع ت}{تحليل}{المتتاليات العددية}{02 سا}{المتتاليات الحسابية}%
{\begin{itemize}\item توليد متتالية ، اتجاه التغير.\end{itemize}}%
{\begin{itemize}\item التعرف على متتالية حسابية وحساب حدها العام.\end{itemize}}%
{المنهاج ، الكتاب المدرسي ، مراجع ، الحاسبة.}
\begin{Sayr}
\Row{مرحلة الانطلاق}{%
\textbf{التهيئة النفسية:} مرتب يبدأ بـ $10000$ DA ويزيد بنفس المبلغ كل سنة: كيف نحسب المرتب بعد $10$ سنوات؟}{10 د}{وضعية مهنية.}
\Row{التشخيص والاكتشاف}{%
\Lab{cEx}{مناقشة نشاط 05 صفحة 179}\par\ck\ المرتب في السنة $n+1$ = المرتب في السنة $n$ + الزيادة الثابتة.\par
\ck\ نستنتج المرتب في السنة $n$ مباشرة بدلالة $n$.\par
}{35 د}{نشاط الكتاب.}
\Row{بناء المعارف}{%
\Lab{cDef}{تعريف}\ $(u_n)$ \textbf{حسابية} إذا وُجد عدد $r$ (الأساس) حيث $u_{n+1}=u_n+r$ من أجل كل $n$.\par\Lab{cThm}{خاصية}\par\ck\ الحد العام: $u_n=u_0+nr$ ، وعموماً $u_n=u_p+(n-p)r$.\par
\ck\ إذا كان $r>0$ فهي متزايدة ، $r<0$ متناقصة ، $r=0$ ثابتة.\par
}{50 د}{الحد العام يُحفظ.}
\Row{}{%
\Lab{cMeth}{إثبات أنها حسابية}\ نحسب $u_{n+1}-u_n$ ونبين أنه ثابت.\par\Lab{cEx}{مثال}\ $u_3=7$ و $u_8=22$: $5r=15$ ، $r=3$ ، $u_0=7-9=-2$ ، $u_n=3n-2$.\par}{}{}
\Row{}{%
\Lab{cThm!80!black}{حذار}\ $u_{n+1}-u_n$ يجب أن يكون مستقلاً عن $n$.}{}{}
\Row{التقويم}{%
\Lab{cEx}{تطبيق}\par بيّن أن $u_n=5-2n$ حسابية ، عيّن أساسها واتجاه تغيرها ، ثم احسب $u_{20}$.\par\textbf{الحل:} $u_{n+1}-u_n=-2$: حسابية أساسها $-2$ ، متناقصة ؛ $u_{20}=-35$.\par\textbf{واجب منزلي:} تمارين 53 إلى 66 من باب المتتاليات.\par\smallskip\textbf{ملاحظات حول سير الحصة:} \dotfill}{25 د}{تقويم فردي.}
\end{Sayr}
\newpage
\ModhHead{60}{2 ع ت}{تحليل}{المتتاليات العددية}{02 سا}{مجموع حدود متتابعة لمتتالية حسابية ، الوسط الحسابي}%
{\begin{itemize}\item المتتاليات الحسابية.\end{itemize}}%
{\begin{itemize}\item حساب مجموع حدود متتابعة لمتتالية حسابية.\item استعمال الوسط الحسابي.\end{itemize}}%
{المنهاج ، الكتاب المدرسي ، مراجع ، الحاسبة.}
\begin{Sayr}
\Row{مرحلة الانطلاق}{%
\textbf{التهيئة النفسية:} كيف حسب غاوس وهو تلميذ صغير $1+2+\cdots+100$ في دقيقة؟}{10 د}{قصة غاوس.}
\Row{التشخيص والاكتشاف}{%
\Lab{cEx}{وضعية مقترحة}\par\ck\ نكتب المجموع مرتين: بالترتيب وبعكسه ، ثم نجمع: $2S=100\times101$ ، $S=5050$.\par
}{35 د}{الكتابة مرتين.}
\Row{بناء المعارف}{%
\Lab{cThm}{خاصية}\ مجموع حدود متتابعة لمتتالية حسابية: $S=N\times\dfrac{a+b}2$ حيث $N$ عدد الحدود ، $a$ الحد الأول و $b$ الحد الأخير.\par\ck\ عدد الحدود من $u_p$ إلى $u_n$ هو $n-p+1$.\par
\Lab{cEx}{مثال}\ $u_n=3n-2$: $u_0+u_1+\cdots+u_{10}=11\times\dfrac{-2+28}2=143$.\par}{50 د}{الدستور يُحفظ بالكلمات.}
\Row{}{%
\Lab{cThm}{الوسط الحسابي}\ $a$ ، $b$ ، $c$ حدود متتابعة لمتتالية حسابية إذا وفقط إذا كان $2b=a+c$.\par\Lab{cThm!80!black}{حذار}\ نعدّ الحدود بدقة: من $u_0$ إلى $u_{10}$ هناك $11$ حداً.}{}{}
\Row{التقويم}{%
\Lab{cEx}{تطبيق}\par احسب $S=3+7+11+\cdots+79$.\par\textbf{الحل:} حسابية أساسها $4$ ؛ عدد الحدود $\dfrac{79-3}4+1=20$ ؛ $S=20\times\dfrac{3+79}2=820$.\par\textbf{واجب منزلي:} تمارين 69 ، 70 ، 72 ، 73 من باب المتتاليات.\par\smallskip\textbf{ملاحظات حول سير الحصة:} \dotfill}{25 د}{تقويم فردي.}
\end{Sayr}
\newpage
\ModhHead{61}{2 ع ت}{تحليل}{المتتاليات العددية}{02 سا}{المتتاليات الهندسية}%
{\begin{itemize}\item المتتاليات الحسابية ، النسب المئوية.\end{itemize}}%
{\begin{itemize}\item التعرف على متتالية هندسية وحساب حدها العام.\item نمذجة ظاهرة (فوائد مركبة ، نمو).\end{itemize}}%
{المنهاج ، الكتاب المدرسي ، مراجع ، الحاسبة.}
\begin{Sayr}
\Row{مرحلة الانطلاق}{%
\textbf{التهيئة النفسية:} مبلغ $50000$ DA موضوع في بنك بفائدة مركبة $4\,\%$ سنوياً: كم يصبح بعد سنتين؟}{10 د}{وضعية بنكية.}
\Row{التشخيص والاكتشاف}{%
\Lab{cEx}{وضعية مقترحة}\par\ck\ $C_1=50000\times1{,}04=52000$ ، $C_2=52000\times1{,}04=54080$.\par
\ck\ نمر من حد إلى الموالي بالضرب في نفس العدد: $C_{n+1}=1{,}04\,C_n$.\par
}{35 د}{الحاسبة.}
\Row{بناء المعارف}{%
\Lab{cDef}{تعريف}\ $(v_n)$ \textbf{هندسية} إذا وُجد عدد $q$ (الأساس) حيث $v_{n+1}=q\,v_n$ من أجل كل $n$.\par\Lab{cThm}{خاصية}\par\ck\ الحد العام: $v_n=v_0\,q^n$ ، وعموماً $v_n=v_p\,q^{n-p}$.\par
\ck\ إذا كان $v_0>0$: $q>1$ متزايدة ، $0<q<1$ متناقصة.\par
}{50 د}{نقارن مع الحسابية في جدول.}
\Row{}{%
\Lab{cMeth}{إثبات أنها هندسية}\ نكتب $v_{n+1}$ على الشكل $q\times v_n$ مع $q$ ثابت.\par\Lab{cEx}{مثال}\ $v_2=12$ و $v_5=96$: $q^3=8$ ، $q=2$ ، $v_0=3$ ، $v_n=3\times2^n$ ؛ $C_n=50000\times1{,}04^n$.\par}{}{}
\Row{}{%
\Lab{cThm!80!black}{حذار}\ إذا كان $q<0$ فالمتتالية ليست رتيبة.}{}{}
\Row{التقويم}{%
\Lab{cEx}{تطبيق}\par بيّن أن $v_n=2\times\left(\frac13\right)^n$ هندسية ، وعيّن اتجاه تغيرها.\par\textbf{الحل:} $v_{n+1}=\frac13v_n$: هندسية أساسها $\frac13$ وحدها الأول $2>0$ ، إذن متناقصة.\par\textbf{واجب منزلي:} تمارين 74 ، 76 ، 77 ، 79 ، 81 ، 82 من باب المتتاليات.\par\smallskip\textbf{ملاحظات حول سير الحصة:} \dotfill}{25 د}{تقويم فردي.}
\end{Sayr}
\newpage
\ModhHead{62}{2 ع ت}{تحليل}{المتتاليات العددية}{02 سا}{مجموع حدود متتابعة لمتتالية هندسية ، الوسط الهندسي}%
{\begin{itemize}\item المتتاليات الهندسية.\end{itemize}}%
{\begin{itemize}\item حساب مجموع حدود متتابعة لمتتالية هندسية.\item استعمال الوسط الهندسي.\end{itemize}}%
{المنهاج ، الكتاب المدرسي ، مراجع ، الحاسبة.}
\begin{Sayr}
\Row{مرحلة الانطلاق}{%
\textbf{التهيئة النفسية:} أسطورة رقعة الشطرنج: حبة قمح في الخانة الأولى ، ثم نضاعف في كل خانة: كم حبة في المجموع؟}{10 د}{الأسطورة.}
\Row{التشخيص والاكتشاف}{%
\Lab{cEx}{وضعية مقترحة}\ $S=1+2+2^2+\cdots+2^{10}$.\par\ck\ $2S-S=2^{11}-1$ ، إذن $S=2047$.\par
}{35 د}{طريقة $qS-S$.}
\Row{بناء المعارف}{%
\Lab{cThm}{خاصية}\ إذا كان $q\neq1$: $S=a\times\dfrac{1-q^{N}}{1-q}$ حيث $a$ الحد الأول و $N$ عدد الحدود.\par\Lab{cEx}{مثال}\ $1+2+\cdots+2^{63}=2^{64}-1$ (عدد هائل!).\par}{50 د}{مقارنة الدستورين.}
\Row{}{%
\Lab{cThm}{الوسط الهندسي}\ $a$ ، $b$ ، $c$ حدود متتابعة لمتتالية هندسية إذا وفقط إذا كان $b^2=ac$.\par\Lab{cThm!80!black}{حذار}\ $N$ هو عدد الحدود وليس دليل الحد الأخير.}{}{}
\Row{التقويم}{%
\Lab{cEx}{تطبيق}\par احسب $S=3+6+12+\cdots+3\times2^7$.\par\textbf{الحل:} الحد الأول $3$ ، $q=2$ ، $N=8$: $S=3\times\dfrac{1-2^8}{1-2}=765$.\par\textbf{واجب منزلي:} تمارين 83 ، 84 ، 85 ، 88 من باب المتتاليات.\par\smallskip\textbf{ملاحظات حول سير الحصة:} \dotfill}{25 د}{تقويم فردي.}
\end{Sayr}
\newpage
\ModhHead{63}{2 ع ت}{تحليل}{المتتاليات العددية}{02 سا}{نهاية متتالية ، المتتاليات المتقاربة}%
{\begin{itemize}\item نهايات الدوال عند $+\infty$ ، المتتاليات الحسابية والهندسية.\end{itemize}}%
{\begin{itemize}\item حساب نهاية متتالية ، والتعرف على متتالية متقاربة.\end{itemize}}%
{المنهاج ، الكتاب المدرسي ، مراجع ، الحاسبة.}
\begin{Sayr}
\Row{مرحلة الانطلاق}{%
\textbf{التهيئة النفسية:} $\left(\frac12\right)^n$ لما يكبر $n$: ماذا يحدث؟ و $2^n$ ؟}{10 د}{الحاسبة.}
\Row{التشخيص والاكتشاف}{%
\Lab{cEx}{وضعية مقترحة}\par\ck\ جدول قيم بالحاسبة لـ $0{,}5^n$ و $2^n$ و $\dfrac{2n+1}{n+3}$.\par
\ck\ نلاحظ: الأولى تقترب من $0$ ، الثانية تكبر بلا حدود ، الثالثة تقترب من $2$.\par
}{35 د}{جداول قيم.}
\Row{بناء المعارف}{%
\Lab{cDef}{تعريف}\ $(u_n)$ \textbf{متقاربة} إذا كانت لها نهاية حقيقية $l$ لما $n\to+\infty$ ، وإلا فهي متباعدة.\par\Lab{cThm}{خاصية}\par\ck\ إذا كان $u_n=f(n)$ و $\lim\limits_{x\to+\infty}f(x)=l$ فإن $\lim u_n=l$.\par
\ck\ $q^n$: نهايتها $0$ إذا كان $-1<q<1$ ، و $+\infty$ إذا كان $q>1$.\par
\ck\ متتالية حسابية أساسها $r>0$ نهايتها $+\infty$.\par
}{50 د}{ملخص النهايات المرجعية.}
\Row{}{%
\Lab{cEx}{مثال}\par\ck\ $\lim\dfrac{2n+1}{n+3}=2$: متقاربة.\par
\ck\ $S_n=1+\frac12+\cdots+\left(\frac12\right)^n=2\left(1-\left(\frac12\right)^{n+1}\right)$ ، $\lim S_n=2$.\par
\Lab{cThm!80!black}{حذار}\ $\lim(-1)^n$ غير موجودة: المتتالية متباعدة رغم أنها لا تؤول إلى $\infty$.}{}{}
\Row{التقويم}{%
\Lab{cEx}{تطبيق}\par ادرس تقارب: $w_n=5-2\times0{,}8^n$ ؛ $t_n=n^2-3n$.\par\textbf{الحل:} $0{,}8^n\to0$: $w_n\to5$ ، متقاربة ؛ $t_n\to+\infty$: متباعدة.\par\textbf{واجب منزلي:} تمارين 51 ، 52 ، 71 ، 82 ، 87 من باب المتتاليات.\par\smallskip\textbf{ملاحظات حول سير الحصة:} \dotfill}{25 د}{نهاية المحور: فرض أو مراجعة.}
\end{Sayr}
\end{document}
